\documentclass[openany,fontset=fandol]{ctexbook}
\usepackage{amsmath,amssymb,amsthm}
\usepackage[unicode,hidelinks]{hyperref}
\setlength{\emergencystretch}{2em}
\raggedbottom
\DeclareMathOperator{\Spec}{Spec}
\DeclareMathOperator{\pr}{pr}
\begin{document}
\title{代数几何基础第五卷预备笔记\\法文覆盖审计：任意超平面截面的连通性}
\author{第96批\quad 法、英、中文三方校订}
\date{2026年11月23日}
\maketitle

\section*{范围与接续}
本批对读法文物理24页下部、25页及英译物理18页：第6节引言、命题6.1、证明和作者的编写建议。法25页页末划销的下一节标题另作校记。连续法校至物理25页末。法26页上部另有一段划销旧稿，其后接第7节；两者下一批核。第三辑1158页合订本未改动。

\section*{6\quad 任意超平面截面的连通性}
这里将泛超平面截面的几何连通性判据（3.3）与 Zariski 连通性定理结合，得到任意超平面截面的连通性结论。

\paragraph{命题 6.1.}
设 $S=\Spec k$，$k$ 为代数闭域，$X$ 在 $k$ 上固有。\footnote{法24页右缘在域的假设之后手写补入 $X$ 的固有性；法25页证明又明确以 $X$ 在 $k$ 上固有为前提。英译在正文方括号保留这一条件，脚注称字迹不清。这里将它列入正文，不视为可选条件。}
设对 $X$ 的每个不可约分支 $X_i$，均有 $\dim\overline{f(X_i)}\geq2$。再设不存在闭子集 $Z\subset X$，使得 $\dim\overline{f(Z)}\leq0$ 且 $X\setminus Z$ 不连通。在这些条件下，对每个 $\xi\in P^\vee$，纤维 $Y_\xi$ 都在 $k(\xi)$ 上几何连通。\footnote{法24末以 $f(X_i)$ 而非 $f(X)$ 为像；法25首行的排除对象是闭子集 $Z$，不是 $X$ 的连通分支。空集也纳入这里的删集条件。此处 $f(X_i)$ 以及 $f(Z)$ 在固有假设下都已闭，写闭包不改变条件。}

\begin{proof}
由于每个 $f(X_i)$ 都不是有限集，$Y$ 的每个不可约分支都占优于 $P^\vee$。另一方面，$Y\to P^\vee$ 是固有态射：$Y\to X$ 是固有的，$X$ 在 $k$ 上固有，而 $P^\vee$ 在 $k$ 上分离。按（3.3），$Y\to P^\vee$ 的泛纤维 $Y_\eta$ 几何连通。最后，$P^\vee$ 正则，因而几何单支。应用正式出版的 EGA~IV 第三分册命题（15.5.3），即 Zariski 连通性定理的一种推论，可知 $Y\to P^\vee$ 的所有纤维都几何连通。\footnote{法25及英18均引作（15.6.3）。正式 EGA~IV\textsubscript{3} 印刷页237的（15.6.3）讨论沿截面的纤维连通分支与普遍开性，并非这里所需断言；印刷页233的（15.5.3）恰给出本证明需要的下半连续性。本文按正式出版编号订正，不断定手稿错号产生的原因。}
\end{proof}

\paragraph{作者给编稿者的说明.}
同类证明还可沿（4.4）的方向推广（6.1）；若不打算做这个练习，至少应作为备注提及。也要说明，（6.1）并不排除维数超过预期的超平面截面。从编排上看，把包括（3.3）和（4.4）在内的所有连通性问题集中到同一节，也许更清楚。\footnote{法25的前后两个命题号有手改；依本段位置和英译读为6.1。作者说“可以推广”并未给出精确命题与证明；此处不把它写成已建立的结果。参考本系列第93批对4.4适用范围的校记。}

\section*{译校补证：如何从泛纤维推出全部纤维}
令 $q:Y\to B=P^\vee$。关联族 $Y$ 是 $X\times_k B$ 内的闭子概形，所以 $q$ 是固有的。投影 $Y\to X$ 在局部是射影空间丛：关联方程描述经过 $f(x)$ 的超平面，在固定一点的仿射坐标片上可消去一个超平面系数。因此，在每个不可约分支 $X_i$ 上的逆像是不可约的，这些逆像给出 $Y$ 的不可约分支。任意 $k$-超平面都与正维闭射影集 $f(X_i)$ 相交。这些逆像对 $B$ 的像是闭集，含全部 $k$-点，故为整个 $B$。特别地，$q$ 满射。

记 $n(\xi)$ 为 $Y_\xi$ 在 $k(\xi)$ 的代数闭包上换底后的连通分支数。正式 EGA~IV\textsubscript{3}（15.5.3）断言：在不可约、局部诺特基底上，若态射固有且源的每个不可约分支都占优于基底，则 $n$ 在几何单支点下半连续。这里每个点都是几何单支点，且泛点 $\eta$ 是每个点的泛化，故 $n(\xi)\leq n(\eta)=1$。又由满射性 $n(\xi)\ge1$，所以 $n(\xi)=1$。这一步并不把“泛纤维几何连通”替换成“泛纤维的全局函数环就是泛点剩余域”。例如非平凡有限纯不可分域扩张的谱几何连通，却不满足后一条件。

出版对照：\href{https://www.numdam.org/item/PMIHES_1966__28__5_0/}{EGA IV第三分册，命题（15.5.3），印刷页233；比较（15.6.3），印刷页237}。正式证明经过 Stein 分解，联系有限态射各纤维的几何点数与原态射的几何连通分支数；亦可对照 \href{https://stacks.math.columbia.edu/tag/03GX}{Stacks Project，§37.53}。

\section*{译校补证：固有性为何不能省略}
第三辑旧核查报告 B2-H09 的反例可直接验证。令
\[
 X=\Spec\bigl(\mathbf C[x,y,z]/(xy-1)\bigr)
   \simeq \mathbf G_m\times\mathbf A^1,
 \qquad
 f:X\hookrightarrow\mathbf A^3\hookrightarrow\mathbf P^3.
\]
这是二维不可约非固有概形。任何像的闭包至多零维的闭子集 $Z$ 都只含有限个点，因为 $f$ 是浸入；去掉这些点仍是非空不可约空间。可是超平面 $x+y=0$ 的逆像具有坐标环
\[
 \mathbf C[x,z]/(x^2+1)\simeq \mathbf C[z]\times\mathbf C[z] .
\]
两条仿射直线是不相交的，故截面不连通。这个例子只反驳删去固有性的版本。

\section*{法25页末的版面校记}
页末下一节标题整行被划销，印号为旧号6；其上方可见手补“或多截面”。可辨题意为“应用于构造指定类型的超平面截面或多截面”。英18页底、英19页顶把它排为第7节，法26页后半也有相应的第7节标题。本批只登记页末标题，不将划销标题计作另一节，不声称已经译校第7.1条。
\end{document}

